\scanpage{59}

\uline{Th. 4} Équivalence de catégories entre

a) Cat. des couples $(X,K,S)$ $X$ surface conforme à bord compacte, $(K,S)$ complexe top. \inserted{cpact} \deleted{\unsure{et|de|ext}} avec $\partial X\subset K\subset X$, $K$ assez régulier (rectifiable suffit)

% slanted
\margin{\uline{NB} On ne suppose pas les \unsure{cptes.|comp.|cpts.} conn. de $K\smallsetminus S$ des intervalles (il peut y avoir des cercles)}

b) Cat des systèmes $(\widetilde X,\widetilde S,\sigma)$, où $\widetilde X$ est une surface conforme à bord \inserted{compacte}, $\widetilde S$ une partie finie de $\partial\widetilde X$, $\sigma$ une involution de $\operatorname{Déc}(\partial\widetilde X/\widetilde S)$, $\sigma$ assez régulière (rectifiable suffit).

Le foncteur a) $\longrightarrow$ b) résulte du th. \unsure{3|1}, avec un complément qui explicite les propriétés de régularité de la fonction-recollement dont il a été question dans le lemme 2. Pour le foncteur en sens inverse, il est connu topologiquement, et comme faisceau structural conforme (large) sur $X$\inserted{$=\widetilde X/\unsure{\mathcal R|R}$}, on prend les fonctions continues qui sont conformes dans $X\smallsetminus K\simeq\operatorname{Int}\widetilde X$.
% archive sheet number: 61
